% Swirl paper --- a map, not a proof.
% Ordinary unaugmented Navier--Stokes leftover after CF / BdVB.
% Not a close. Not Fefferman (C)/(D). Not Theorem A retitled.
% Compile with pdflatex when a TeX tree is present.
% Phone PDF: python3 scripts/da_swirl_pdf.py
\documentclass[11pt]{article}
\usepackage[margin=1in]{geometry}
\usepackage{amsmath,amssymb,amsthm}
\usepackage[T1]{fontenc}
\usepackage{microtype}
\usepackage{hyperref}

\title{Swirl, alignment, and the leftover in\\
three-dimensional Navier--Stokes\\
\large A map, not a proof}
\author{Jonathan Robert Simons\\
Prime Field Technologies}
\date{10 September 2026}

\begin{document}
\maketitle

\begin{abstract}
The geometric swirl leftover after Constantin--Fefferman (1993)
and Beir\~ao da Veiga--Berselli (2002) is recorded. Alignment of
vorticity direction at Lipschitz scale, and then at H\"older
exponent \(1/2\), is a theorem as an \emph{if}. The pairs those
papers refused remain. This note names the missing local bound
WRITE~(6) and does not prove it. A forced finite-time construction
with a smooth external force, announced 8~September 2026, addresses
a different official statement (Fefferman (C)/(D)) from unforced
regularity ((A)/(B)). It is not WRITE~(6). A separate note records
global regularity for a Ladyzhenskaya / \(p\)-Laplacian augmentation;
that is a different equation.
\end{abstract}

\section{What swirl is, in this note}

Let \(u\) be a divergence-free velocity in three dimensions,
\(\omega=\nabla\times u\), \(\xi=\omega/|\omega|\) on
\(\{\omega\neq 0\}\), and \(\varphi(x,y)\) the angle between
\(\xi(x)\) and \(\xi(y)\). Vortex stretching is
\[
\alpha(x)=\mathrm{P.V.}\int D(\hat z,\xi(x),\xi(y))
\frac{|\omega(y)|}{|x-y|^3}\,dy,
\qquad
|D|\le C|\sin\varphi|.
\]
The local enstrophy identity on a cutoff \(\phi\) supported in a
ball of radius \(r\) has stretching \(\int\alpha|\omega|^2\phi\)
on the right and dissipation \(\nu\int|\nabla\omega|^2\phi\) on
the left. This is the classical unaugmented equation.

\section{The cut that sits, as an if}

Constantin--Fefferman (1993): if
\(|\sin\varphi(x,y)|\le|x-y|/\rho\) whenever both
\(|\omega|\ge\Lambda\), the solution is strong on \([0,T]\).
Beir\~ao da Veiga--Berselli (2002): H\"older \(1/2\) suffices.
Alignment drops the kernel. Call those pairs Good. The complementary
intense pairs with \(|\sin\varphi|>C_*|x-y|^{1/2}\) are Bad.
Good pairs are their theorem. Bad pairs are the pairs they refused.
Gruji\'c (2009) localizes the \emph{condition}, not the bad-pair
integral.

\section{WRITE (6)}

On a parabolic cylinder \(Q_r\), Bad pairs only,
\[
A_{\mathrm{bad}}(Q_r)
\le
\frac\nu8\iint_{Q_r}|\nabla\omega|^2\phi
+C r^{-2}\iint_{Q_r}|\omega|^2.
\]
The kernel on Bad is still \(|z|^{-3}\). Hardy--Littlewood--Sobolev
returns local \(E^3\). The estimate is stated. It is not proved.

\section{A cylinder still needs more than (6)}

Skin flux at energy level is Caffarelli--Kohn--Nirenberg as
smallness, not as an a priori from finite energy. Exterior
Biot--Savart is not absorbed as \(r\to 0\) from energy. Local
Serrin after a closed enstrophy budget is the right implication;
it is not a substitute for (6). If (6) sits and the skin flux
from energy does not, the cylinder is still open.

\section{A forced vortex is not this leftover}

Fefferman's official statement: (A)/(B) unforced regularity;
(C)/(D) breakdown with a smooth force allowed. A finite-time
singularity with a smooth external force, announced
8~September 2026, lives in (C) or (D) if it holds. It does not
decide (A) or (B) and does not bound WRITE~(6). This author did
not write that construction and does not claim it.

\section{What is claimed}

The map above. A separate write-up of global regularity for
Navier--Stokes plus Ladyzhenskaya / \(p\)-Laplacian stress at
\(\varepsilon>0\), \(\beta\ge 1/2\) (this PDE only). Inverse-GCD
facts posted as Zenodo 22045478, not as the Riemann hypothesis.
Not claimed: ordinary unforced Navier--Stokes; WRITE~(6) as a
theorem; axisymmetric ordinary NS with swirl as finished; the
forced construction of 8~September 2026.

\begin{thebibliography}{11}
\bibitem{L34} J.~Leray, \textit{Acta Math.} 63 (1934).
\bibitem{S62} J.~Serrin, \textit{Arch.\ Rational Mech.\ Anal.} 9 (1962).
\bibitem{CKN} L.~Caffarelli, R.~Kohn, L.~Nirenberg,
\textit{Comm.\ Pure Appl.\ Math.} 35 (1982), 771--831.
\bibitem{CF} P.~Constantin, C.~Fefferman,
\textit{Indiana Univ.\ Math.\ J.} 42 (1993), 775--789.
\bibitem{BdVB} H.~Beir\~ao da Veiga, L.~C.~Berselli,
\textit{Diff.\ Int.\ Eq.} 15 (2002), 345--356.
\bibitem{NRS} J.~Ne\v{c}as, M.~R\r{u}\v{z}i\v{c}ka, V.~\v{S}ver\'ak,
\textit{Acta Math.} 176 (1996).
\bibitem{ESS} L.~Escauriaza, G.~Seregin, V.~\v{S}ver\'ak,
\textit{Uspekhi Mat.\ Nauk} 58 (2003).
\bibitem{G09} Z.~Gruji\'c, \textit{Comm.\ Math.\ Phys.} 290 (2009).
\bibitem{Feff} C.~L.~Fefferman,
\textit{Existence and smoothness of the Navier--Stokes equation},
Clay Mathematics Institute.
\bibitem{L69} O.~A.~Ladyzhenskaya,
\textit{The Mathematical Theory of Viscous Incompressible Flow}, 1969.
\bibitem{MNR} J.~M\'alek, J.~Ne\v{c}as, M.~R\r{u}\v{z}i\v{c}ka,
\textit{Weak and Measure-valued Solutions to Evolutionary PDEs}, 1996.
\end{thebibliography}

\end{document}
